% Options for packages loaded elsewhere
\PassOptionsToPackage{unicode}{hyperref}
\PassOptionsToPackage{hyphens}{url}
%
\documentclass[
]{book}
\usepackage{amsmath,amssymb}
\usepackage{iftex}
\ifPDFTeX
  \usepackage[T1]{fontenc}
  \usepackage[utf8]{inputenc}
  \usepackage{textcomp} % provide euro and other symbols
\else % if luatex or xetex
  \usepackage{unicode-math} % this also loads fontspec
  \defaultfontfeatures{Scale=MatchLowercase}
  \defaultfontfeatures[\rmfamily]{Ligatures=TeX,Scale=1}
\fi
\usepackage{lmodern}
\ifPDFTeX\else
  % xetex/luatex font selection
\fi
% Use upquote if available, for straight quotes in verbatim environments
\IfFileExists{upquote.sty}{\usepackage{upquote}}{}
\IfFileExists{microtype.sty}{% use microtype if available
  \usepackage[]{microtype}
  \UseMicrotypeSet[protrusion]{basicmath} % disable protrusion for tt fonts
}{}
\makeatletter
\@ifundefined{KOMAClassName}{% if non-KOMA class
  \IfFileExists{parskip.sty}{%
    \usepackage{parskip}
  }{% else
    \setlength{\parindent}{0pt}
    \setlength{\parskip}{6pt plus 2pt minus 1pt}}
}{% if KOMA class
  \KOMAoptions{parskip=half}}
\makeatother
\usepackage{xcolor}
\usepackage{longtable,booktabs,array}
\usepackage{calc} % for calculating minipage widths
% Correct order of tables after \paragraph or \subparagraph
\usepackage{etoolbox}
\makeatletter
\patchcmd\longtable{\par}{\if@noskipsec\mbox{}\fi\par}{}{}
\makeatother
% Allow footnotes in longtable head/foot
\IfFileExists{footnotehyper.sty}{\usepackage{footnotehyper}}{\usepackage{footnote}}
\makesavenoteenv{longtable}
\usepackage{graphicx}
\makeatletter
\def\maxwidth{\ifdim\Gin@nat@width>\linewidth\linewidth\else\Gin@nat@width\fi}
\def\maxheight{\ifdim\Gin@nat@height>\textheight\textheight\else\Gin@nat@height\fi}
\makeatother
% Scale images if necessary, so that they will not overflow the page
% margins by default, and it is still possible to overwrite the defaults
% using explicit options in \includegraphics[width, height, ...]{}
\setkeys{Gin}{width=\maxwidth,height=\maxheight,keepaspectratio}
% Set default figure placement to htbp
\makeatletter
\def\fps@figure{htbp}
\makeatother
\setlength{\emergencystretch}{3em} % prevent overfull lines
\providecommand{\tightlist}{%
  \setlength{\itemsep}{0pt}\setlength{\parskip}{0pt}}
\setcounter{secnumdepth}{5}
\usepackage{booktabs}
\usepackage{amsthm}
\makeatletter
\def\thm@space@setup{%
  \thm@preskip=8pt plus 2pt minus 4pt
  \thm@postskip=\thm@preskip
}
\makeatother
\ifLuaTeX
  \usepackage{selnolig}  % disable illegal ligatures
\fi
\usepackage[]{natbib}
\bibliographystyle{apalike}
\usepackage{bookmark}
\IfFileExists{xurl.sty}{\usepackage{xurl}}{} % add URL line breaks if available
\urlstyle{same}
\hypersetup{
  pdftitle={Jehill's Investing Diary},
  pdfauthor={Jehill Parikh},
  hidelinks,
  pdfcreator={LaTeX via pandoc}}

\title{Jehill's Investing Diary}
\author{Jehill Parikh}
\date{2025-03-05}

\begin{document}
\maketitle

{
\setcounter{tocdepth}{1}
\tableofcontents
}
\chapter{Welcome, My Personal Investing Diary}\label{welcome-my-personal-investing-diary}

Investing is a journey, and every investor has a story. With strong belief, in learning from experiences, tracking insights, and making informed financial decisions. Whether you're a seasoned investor or just starting out, we hope this blog adds values to personal investments journey, by providing---a place of strategies, knowledge and experenice markets, and refine decision-making.

From market trends to portfolio strategies, so specific examples we cover key investment principles to help you navigate the financial world with confidence, as we explore the ups and downs of investing, share valuable lessons, and grow together toward financial independence.

It is important to note that content on this blog is not a recommendation and no investment decisions or recommendations are provide in any of the content on this blog, and investing inovles risk, and the reader is encourage to consulate a professional to seek advice. The content here is education and not personal recommendation.

Stay curious. Stay disciplined. Invest smarter.

\chapter{Investment Checklist}\label{investment-checklist}

There are many invesing thoeries and concepts, however, what generally is most helpful is need to determine a checklist summary various analysis. At the onset we start with a framework to assess macroeconomic, sector-specific, market-related, and investment-specific factors before making informed investment decisions.

\section{1. Macroeconomic Factors}\label{macroeconomic-factors}

\begin{itemize}
\tightlist
\item
  \textbf{GDP Growth}: Probability of economic expansion or contraction affecting investment opportunities.\\
\item
  \textbf{Interest Rates}: Probability of fluctuations impacting borrowing costs and investment decisions.\\
\item
  \textbf{Inflation}: Probability of moderate inflation benefiting asset values or high inflation increasing costs.\\
\item
  \textbf{Unemployment Rate}: Probability of lower unemployment increasing consumer spending or higher unemployment reducing demand.\\
\item
  \textbf{Government Policies}: Probability of supportive or restrictive policies impacting investment sectors.\\
\item
  \textbf{Currency Movements}: Probability of depreciation increasing import costs or appreciation reducing expenses.\\
\item
  \textbf{Commodity Prices}: Probability of price increases raising costs or decreases lowering expenses.
\end{itemize}

\section{2. Sector-Specific Factors}\label{sector-specific-factors}

\begin{itemize}
\tightlist
\item
  \textbf{Sector Growth Rate}: Probability of sector expansion due to market trends or contraction from adverse conditions.\\
\item
  \textbf{Technological Advancements}: Probability of innovation driving growth or obsolescence impacting competitiveness.\\
\item
  \textbf{Regulations}: Probability of favorable policies supporting growth or restrictive regulations increasing compliance costs.\\
\item
  \textbf{Supply Chain Issues}: Probability of disruptions delaying projects or improvements reducing costs.\\
\item
  \textbf{Demand Trends}: Probability of rising demand boosting investments or declining demand reducing opportunities.\\
\item
  \textbf{Sector Sentiment}: Probability of positive investor sentiment driving growth or negative sentiment limiting investments.
\end{itemize}

\section{3. Market Factors}\label{market-factors}

\begin{itemize}
\tightlist
\item
  \textbf{Market Liquidity}: Probability of high liquidity enabling smooth transactions or low liquidity increasing volatility.\\
\item
  \textbf{Investor Sentiment}: Probability of optimism boosting investment inflows or pessimism leading to sell-offs.\\
\item
  \textbf{Institutional Buying/Selling}: Probability of large-scale institutional investments or withdrawals influencing market trends.\\
\item
  \textbf{Volatility}: Probability of stable market conditions or high volatility affecting investment stability.
\end{itemize}

\section{4. Standard Investment Factors}\label{standard-investment-factors}

\begin{itemize}
\tightlist
\item
  \textbf{Value (P/E)}: Probability of undervaluation presenting a buying opportunity or overvaluation limiting returns.\\
\item
  \textbf{Quality (ROE)}: Probability of strong returns on equity driving investment attractiveness or weak returns impacting valuation.\\
\item
  \textbf{Momentum (Price Trend)}: Probability of upward momentum indicating growth potential or downward trends signaling caution.\\
\item
  \textbf{ROCE (Return on Capital Employed)}: Probability of efficient capital use generating high returns or inefficiencies lowering profitability.\\
\item
  \textbf{Size (Market Capitalization)}: Probability of a company or sector gaining large-cap status or remaining in small/mid-cap range.\\
\item
  \textbf{Volatility (Beta)}: Probability of low beta indicating stability or high beta signaling sensitivity to market conditions.\\
\item
  \textbf{Alpha (Risk-Adjusted Return)}: Probability of strong risk-adjusted performance or underperformance relative to benchmarks.\\
\item
  \textbf{Beta (Systematic Risk)}: Probability of low systematic risk for stable investments or high beta indicating exposure to market fluctuations.
\end{itemize}

\section{5. Growth Investing Factors}\label{growth-investing-factors}

\begin{itemize}
\tightlist
\item
  \textbf{Growth Rate}: Probability of sustained revenue and earnings growth over time.\\
\item
  \textbf{Quality of Growth}: Probability of growth being driven by strong fundamentals versus temporary market conditions.\\
\item
  \textbf{Longevity of Growth}: Probability of a company's or sector's growth being sustainable in the long run.\\
\item
  \textbf{Scalability}: Probability of a business model being able to scale efficiently without significant cost increases.\\
\item
  \textbf{Market Expansion}: Probability of entering new markets or expanding market share contributing to sustained growth.\\
\item
  \textbf{Innovation \& R\&D}: Probability of continuous innovation driving competitive advantage and future revenue streams.\\
\item
  \textbf{Competitive Positioning}: Probability of maintaining or strengthening market leadership against competitors.\\
\item
  \textbf{Profitability of Growth}: Probability of revenue growth translating into actual profitability rather than just higher costs.
\end{itemize}

\section{6. Balance Sheet Factors}\label{balance-sheet-factors}

\begin{itemize}
\tightlist
\item
  \textbf{Debt Levels}: Probability of manageable debt improving financial stability versus excessive debt increasing financial risk.\\
\item
  \textbf{Liquidity Ratios}: Probability of strong liquidity ensuring smooth operations versus liquidity constraints limiting flexibility.\\
\item
  \textbf{Asset Quality}: Probability of high-quality assets appreciating in value versus low-quality assets leading to impairments.\\
\item
  \textbf{Cash Flow Stability}: Probability of consistent positive cash flow supporting growth versus erratic cash flow affecting financial planning.\\
\item
  \textbf{Capital Structure}: Probability of an optimal mix of debt and equity enhancing financial resilience versus an unbalanced structure increasing risk.\\
\item
  \textbf{Working Capital Efficiency}: Probability of efficient working capital management improving operations versus inefficiencies leading to cash flow problems.\\
\item
  \textbf{Dividend Sustainability}: Probability of consistent and growing dividends indicating financial strength versus unsustainable payouts leading to cuts.\\
\item
  \textbf{Book Value Growth}: Probability of rising book value per share reflecting increasing intrinsic worth versus declining book value signaling potential issues.
\end{itemize}

\section{7. Operating Leverage Factors}\label{operating-leverage-factors}

\begin{itemize}
\tightlist
\item
  \textbf{Fixed vs.~Variable Costs}: Probability of a high fixed-cost structure amplifying profit growth versus variable costs providing flexibility.\\
\item
  \textbf{Break-even Point}: Probability of reaching the break-even point faster with high operating leverage versus slower with low leverage.\\
\item
  \textbf{Margin Expansion}: Probability of improving operating margins as revenue grows due to fixed cost absorption.\\
\item
  \textbf{Cost Efficiency}: Probability of leveraging economies of scale to reduce per-unit costs and improve profitability.\\
\item
  \textbf{Revenue Sensitivity}: Probability of revenue fluctuations significantly impacting profitability due to high operating leverage.\\
\item
  \textbf{Operational Flexibility}: Probability of adapting cost structures to changing market conditions efficiently.\\
  ```
\end{itemize}

  \bibliography{book.bib,packages.bib}

\end{document}
